%%
%% winter-lecture24.tex
%% 
%% Made by Alex Nelson <pqnelson@gmail.com>
%% Login   <alex@lisp>
%% 
%% Started on  2026-02-28T07:49:02-0800
%% Last update 2026-02-28T07:49:02-0800
%% 

\lecture{}

\begin{lemma}\label{lemma:17.6}
If $A$ is an integral domain and $B$ is an integral extension of $A$,
then $A$ is a field if and only if $B$ is a field.
\end{lemma}

\begin{proof}
\forwardproof\ Assume $A$ is a field. Every $\alpha\in B$ is
integral over $A$ (since $B$ is an integral extension of $A$). So then
there is a monic polynomial for which $\alpha$ is a root
\begin{equation}
\alpha^{n}+a_{n-1}\alpha^{n-1}+\dots+a_{0}=0
\end{equation}
with the coefficients $a_{i}\in A$. Then
\begin{equation}
\alpha(\alpha^{n-1}+\dots+a_{1})=-a_{0},
\end{equation}
so
\begin{equation}
\alpha^{-1}=\frac{-1}{a_{0}}(\alpha^{n-1}+\dots+a_{1})\in A.
\end{equation}
Hence $\alpha^{-1}\in B$, so $B$ must be a field.

\backwardproof\ Assume now that $B$ is a field. Let $\alpha\in A$
be arbitrary. We will prove $\alpha^{-1}\in A$. By similar reasoning,
we see $\alpha^{-1}\in B$ is algebraic over $A$, so then there is a
monic polynomial $P\in A[x]$ for which $P(\alpha^{-1})=0$. Similar
reasoning as before gives us $\alpha^{-1}\in A$.
\end{proof}

\begin{corollary}\label{cor:17.7}
Let $A$ be an integral domain and let $B$ be an integral extension of $A$.
If $P\in\Spec(A)$ is a prime ideal of $A$ and if $Q\in\Spec(B)$ is a
prime ideal of $B$ lying over $P$ (i.e., $Q\cap A=P$), then $P$ is
maximal if and only if $Q$ is maximal.
\end{corollary}

\begin{proof}
Apply lemmas~\ref{lemma:17.5}, \ref{lemma:17.6} to $A_{P}$ and $B_{Q}$.
\end{proof}

\subsection{Dedekind Domains}

\begin{convention}
In this section discussing Dedekind domains, we will always have $A$
be an integral domain, and $K=\Frac(A)$ its field of fractions.
\end{convention}

\begin{definition}
We say $A$ is a \define{Dedekind Domain} if it is Noetherian,
integrally closed, and every prime ideal is maximal.
\end{definition}

\begin{definition}
A nonzero $A$-module $\mathfrak{A}\subset K$ is called a
\define{Fractional Ideal} of $A$ if there exists an $a\in A$ such that
$a\mathfrak{A}\subset A$ (in particular, $a\mathfrak{A}$ is an ideal
of $A$).
\end{definition}

\begin{lemma}\label{lemma:18.1}
Every finite $A$-submodule $\mathfrak{A}$ of $K$ is a fractional ideal
of $A$.

If $A$ is Noetherian, then the converse also holds.
\end{lemma}

\begin{proof}
Since $\mathfrak{A}$ is a finite $A$-module, it has a finite basis
$\{\alpha_{1},\dots,\alpha_{n}\}\subset K$. By clearing the
denominators of $\alpha_{i}$, there exists an $a\in A$ [namely, the
  product of all the denominators of the $\alpha_{i}$] such that
$a\mathfrak{A}\subset A$. Hence $\mathfrak{A}$ is a fractional ideal
of $A$.

Assume $A$ is Noetherian, assume $\mathfrak{A}$ is a fractional ideal
of $A$. Then, for some $a\in A$,
\begin{equation}
I:=a\mathfrak{A}
\end{equation}
is a finitely-generated ideal of $A$. Then $\mathcal{A}=a^{-1}I$ is a
finite $A$-submodule of $K$.
\end{proof}

\begin{lemma}\label{lemma:18.2}
For a Noetherian integral domain $A$ and a nonzero ideal $I\ideal A$,
there exists prime ideals $P_{1}$, \dots, $P_{n}$ such that their
product $P_{1}(\dots)P_{n}\subset I$.
\end{lemma}

\begin{proof}
Let $\mathcal{I}$ be the set of ideals of $A$ for which the claim does
not hold. Assume for contradiction $\mathcal{I}\neq\emptyset$.
Since $A$ is Noetherian, there exists a maximal element in
$\mathcal{I}$. Let us call such a maximal element
$I\in\mathcal{I}$. Then $I$ cannot be prime. Then there exists $x,y\in A\setminus I$
such that $xy\in I$. Since $I$ is maximal in $\mathcal{I}$, we must
have $I+(x)\notin\mathcal{I}$ and similarly $I+(y)\notin\mathcal{I}$
otherwise $I$ would not be maximal in $\mathcal{I}$ (we would have a
bigger ideal, namely $I+(x)$ or $I+(y)$).
So these ideals must satisfy the claim. Then
\begin{equation}
(I+(x))(I+(y))\in\mathcal{I},
\end{equation}
since $(I+(x))(I+(y))=I+(x)I+(y)I+(x)(y)=I+(xy)\supset P_{1}\dots P_{n}$ 
which is a contradiction (since $I+(x)\supset P_{1}\dots P_{n}$ and
$I+(y)\supset P_{1}\dots P_{n}$ and so $I\supset P_{1}\dots P_{n}$
which contradicts $I\nsupset P_{1}\dots P_{n}$). Hence $\mathcal{I}=\emptyset$.
\end{proof}

\begin{lemma}\label{lemma:18.3}
For a Dedekind domain $A$ and a nonzero prime ideal $P\in\Spec(A)$,
there exists a nonzero element $x\in K\setminus A$ such that
$xP\subset A$.
\end{lemma}

\begin{proof}
For $0\neq a\in P$, by Lemma~\ref{lemma:18.2} there exists prime
ideals $P_{1},\dots,P_{n}\in\Spec(A)$ such that
\begin{equation}
P_{1}\dots P_{n}\subset(a)\subset P.
\end{equation}
Assume that $n$ is taken to be minimal (i.e., any $n-1$ such ideals
fail to satisfy the condition). Since $P$ is prime, by reindexing, we
have $P_{n}\subset P$. Then by definition of a Dedekind domain, all
prime ideals are maximal, so
\begin{equation}
P_{n}=P.
\end{equation}
By minimality of $n$,
\begin{equation}
P_{1}(\dots)P_{n-1}\nsubset P
\end{equation}
So there exists a
\begin{equation}
b\in P_{1}\dots P_{n-1}
\end{equation}
such that $b\notin(a)$. Then
\begin{equation}
x=\frac{b}{a}\notin A.
\end{equation}
Then
\begin{subequations}
  \begin{align}
xP &= a^{-1}bP_{n}\\
&\subset a^{-1}P_{1}(\dots)P_{n-1}P_{n}\\
&\subset a^{-1}(a) = A,
  \end{align}
\end{subequations}
as desired.
\end{proof}

\begin{proposition}\label{prop:18.4}
Suppose $A$ is a Dedekind domain, suppose $\mathfrak{A}$ is a
fractional ideal of $A$. Then there exists another fractional ideal
denoted $\mathfrak{A}^{-1}$ such that $\mathfrak{A}\mathfrak{A}^{-1}=A$.
\end{proposition}

This means the fractional ideals of $A$ form a group when $A$ is a
Dedekind domain. (Also note: by Lemma~\ref{lemma:18.3},
$A\propersubset P^{-1}$ for any prime ideal $P$ of $A$.)

\begin{proof}
Let
\begin{equation}
\mathfrak{A}=\frac{1}{a}I
\end{equation}
for some $a\in A$ and ideal $I\ideal A$. We will show
\begin{equation}
I^{-1}:=\{x\in K\mid xI\subset A\}
\end{equation}
(and by the previous Lemma, $I^{-1}\neq0$), and $I^{-1}$ satisfies the
claim---i.e., if we take $\mathfrak{A}^{-1}=aI^{-1}$. Without loss of
generality, we may take $a=1$. We see $I^{-1}$ is a fractional field
of $I$ [Pf: it is an $A$-submodule of $K$ and $bI^{-1}\subset A$ for
  any $b\in I$, so it must be a fractional ideal of $A$.]

Assume for contradiction $II^{-1}\propersubset A$. Then the ideal
$II^{-1}$ is some maximal ideal $P$ of $A$. Then by
Lemma~\ref{lemma:18.3}, there exists $x\in K\setminus A$ such that
$xII^{-1}\subset xP\subset A$. But $xII^{-1}=(xI^{-1})I$ and
$xI^{-1}\subset I^{-1}$. (For any $y\in xI^{-1}$, we have $yI\subset A$,
so $y\in I^{-1}$.) By Lemma~\ref{lemma:18.1}, $I^{-1}$ isa finite
$A$-submodule of $K$. Then we see $I^{-1}$ has a finite $A$-basis
$\{\alpha_{i}\}$ over $A$. Since $xI^{-1}\subset I^{-1}$, we have
\begin{equation}\label{eq:math120b:lec24:winter2026:char-poly}
x\alpha_{i}=\sum_{j}a_{ij}\alpha_{j}
\end{equation}
for all $i$. This means $x$ satisfies the characteristic polynomial
[Equation~\eqref{eq:math120b:lec24:winter2026:char-poly}] (over $A$),
i.e., $x$ is integral over $A$. But since $A$ is a Dedekind domain
(viz., $A$ is integrally closed) and $x\in K\setminus A$ we get a
contradiction. Therefore $II^{-1}=A$, hence the claim.
\end{proof}

\begin{proposition}\label{prop:18.5}
Let $A$ be a Dedekind domain, let $I\ideal A$ be an ideal of $A$.
Then there exists prime ideals $P_{1}$, \dots, $P_{n}$ of $A$ and
positive integers $e_{1}$, \dots, $e_{n}$ such that
\begin{equation}
I=P_{1}^{e_{1}}(\cdots)P_{n}^{e_{n}}
\end{equation}
and moreover these must be unique.
\end{proposition}

\begin{proof}
\textsc{Claim 1: Existence}.
Suppose that $\mathcal{I}$ is the set of ideals of $A$ with no such
prime factorization. Assume for contradiction $\mathcal{I}\neq\emptyset$.
By the same discussion as in the last proposition, we have
a maximal element $I\in\mathcal{I}$. We have $I\propersubset IP^{-1}\subset A$
where $P\ideal A$ is a maximal ideal containing $I\propersubset P$.
Then $IP^{-1}$ has a prime factorization (since $I\in\mathcal{I}$ is
maximal and $I\propersubset IP^{-1}$, so $IP^{-1}\notin\mathcal{I}$),
and so
\begin{subequations}
  \begin{align}
I &= (IP^{-1})P\\
&=(P_{1}^{e_{1}}(\cdots)P_{n}^{e_{n}})P
  \end{align}
\end{subequations}
is such a prime factorization. Hence contradiction.

\textsc{Claim 2: Uniqueness}. Suppose
\begin{equation}
I=P_{1}(\cdots)P_{m}=Q_{1}(\cdots)Q_{n}
\end{equation}
are two such prime factorizations. Then
\begin{equation}
Q_{1}(\cdots)Q_{n}=I\subset P_{1},
\end{equation}
so by possibly reindexing the $Q_{i}$'s we have
\begin{equation}
Q_{1}\subset P_{1}.
\end{equation}
But since prime ideals are maximal ideals in a Dedekind domain, we
must have
\begin{equation}
Q_{1}=P_{1}.
\end{equation}
Then by applying the same argument inductively, we can reindex to find
$m=n$ and also $Q_{2}=P_{2}$, \dots, $Q_{m}=P_{n}$.
\end{proof}